%%%     Package List and Document Class     %%%%
\documentclass[12pt,twoside]{paper}
\usepackage[numbers,sort&compress]{natbib}
\usepackage{graphics}
\usepackage{graphicx}
\usepackage{makeidx}
\usepackage{amsfonts}
\usepackage{amsmath}
\usepackage{enumerate}
\usepackage{enumitem}
\usepackage[colorlinks,breaklinks]{hyperref}
\usepackage{hyperref}
\usepackage{color}
\usepackage{hypernat}
\usepackage{subfig}
\usepackage{pslatex}
\usepackage{soul}
\usepackage{fancyhdr}
\newcommand{\forceindent}{\leavevmode{\parindent=1em\indent}}
\usepackage{wrapfig}
\usepackage{float}
\usepackage[compact]{titlesec}
\usepackage{mdwlist}
\usepackage[margin=1in]{geometry}
\usepackage{amssymb}
\renewcommand{\headrulewidth}{0pt}
\newcommand{\ti}{Robust Nonintrusive Conservative Limiting for Legacy Codes}
%\newcommand{\piname}{\textcolor{red}{CUI Page Markings. (DELETE or COMMENT OUT if not applicable.)}}
\newcommand{\pinamee}{PI Name (Tokareva, Svtelana)}
\setlength{\headheight}{12pt}
\pagestyle{fancyplain}
\fancyhf{}
\lhead{\fancyplain{}{\scriptsize \ti}}
\rhead{\fancyplain{}{\scriptsize Proposal Number (obtained from the Q system)}}
%\chead{\fancyplain{}{\scriptsize \textcolor{red}{CUI Page Markings (if applicable)}}}
\lfoot{\fancyplain{}{\scriptsize \pinamee}}
\fancyfoot[R]{\thepage}
\renewcommand{\headrulewidth}{0pt}
\setcounter{page}{1}
\setlength{\parindent}{0em}
\setlength{\parskip}{0em}
\raggedright
\raggedbottom
%%%%%%%%%%%%%%%%%%%%%%%%%%%%%%%%%%%%%%%%%%%%%%%%

%%%     Graphics Path     %%%%
\graphicspath{{./fig/}}
%%%%%%%%%%%%%%%%%%%%%%%%%%%%%%%%%%%%%%%%%%%%%%%%
\begin{document}

%%%    Hyperlink Setup     %%%%
\hypersetup{citecolor=blue,
            linkcolor=blue,
            urlcolor=blue,
            pdftitle=ti}
%            pdfauthor=\auth}
%%%%%%%%%%%%%%%%%%%%%%%%%%%%%%%%%%%%%%%%%%%%%%%%







%%%%%%%%%%%%%%%%%%%%%%%%%%%%%%%%%%%%%%%%%%%%%%%
%%%     FEASIBILITY STUDY COVERPAGE     %%%%
%%%%%%%%%%%%%%%%%%%%%%%%%%%%%%%%%%%%%%%%%%%%%%%

\begin{center}
    \large \textbf{Feasibility Study (FS)}\\
    \vspace{0.5cm}
    \normalsize \textbf{Robust Nonintrusive Conservative Limiting for Legacy Codes}\\
    Tokareva, Svetlana; T5; tokareva@lanl.gov
\end{center}
\vspace*{\baselineskip}

\textbf{Duration (mo):} 12 months.
\vspace*{\baselineskip}

\textbf{Funding level (\$K):} \$150K.
\vspace*{\baselineskip}

\textbf{Project Description:}

Sampling and exploring large parameter spaces using LANL's production
codes often fails due to robustness issues. Machine learning
techniques are significantly affected by the robustness issue as they
require numerous simulations to learn the functional dependence on
input parameters.  The failure probabibilty is high when the
models/approximation techniques involve various time/space scales,
multi-rate time stepping, adaptive multi-resolution, and/or mix
approximation paradigms. We posit that ensuring that every piece of
the computational chain maintains physical bounds is the key to
achieving robustness.

One key tool to maintain physical bounds is limiting, but most known
limiting techniques are specialized and intrusive.  They mostly
consist of limiting slopes or fluxes, and are all somewhat based on
the flux-corrected transport paradigm of Boris–Book (1973) and Zalesak
(1979) by which a “non-admissible” but accurate high-order solution is
squeezed in the direction of a “good” low-order solution until it is
meets the requires bounds.  The squeezing idea is powerful,
but this paradigm has four major defects: (i)
It is intrusive: i.e., it requires manipulating slopes and fluxes
which is impossible to do in LANL's legacy codes without engaging in costly
and error-prone restructuration; (ii) It fails when the high-order and
low-order solutions do not carry the same “generalized” mass; (iii) It
cannot be applied to modal approximation techniques and reduced-order
models; (iv)  It fails for heterogeneous
time/space approximations.

The goal of this FS is to develop generic limiting techniques that are
robust, efficient, HPC portable, and potentially easy to implement in
LANL's legacy codes.


\vspace*{\baselineskip}



\textbf{Cover Page Markings:} NA


\pagebreak

\textbf{Problem, Hypothesis \& Technical Rationale:}
\vspace*{\baselineskip}

\textit{Why:} Accurate quantitative predictions of coupled multiphysics  systems
are at the center of LANL's mission and many innovative applications like
hypersonic flight, high energy density phenomena, control and
manipulation of high-speed flows involving complex fluids, mixtures,
and explosives, possibly interacting with radiation and
electromagnetic fields. Some of these
applications have the potential to enable new energy sources by
forming and controlling plasmas with magnetic confinement (tokamaks and
stellarators) or inertial confinement (ICF).  But simulating complex multiphysics
problems involving multiple time and length scales \textbf{still poses}
robustness challenges both for direct and inverse
methods.

Although the machine learning revolution promises to solve some of
these problems, data sampling is a major bottleneck due to \textbf{lack of
robustness} for many sample-driven computational approaches such as
surrogate-based optimization, Bayesian inverse problems, and
uncertainty quantification.  Machine learning requires running legacy
codes a large number of times in order to learn the functional
dependence with respect to input parameters. But sampling large
parameter spaces via thousands of automated runs of complex
multiphysics legacy codes using traditional approximation methods is
hindered by the fact that a significant number of simulations never
reach completion due to robusness issues, thereby undermining the
promises of the machine learning revolution.  There is therefore an
\textbf{urgent need} for the development of agile, non-intrusive, and provably
robust numerical methods that can reliably solve multiphysics
applications using heterogeneous approximation techniques.

\vspace*{\baselineskip}

\textit{What:} We posit that a key to achieving robustness when leveraging
heterogeneous approximation techniques is to ensure that every piece
of the computational chain is maintains physical bounds (i.e. is
invariant-domain preserving (IDP)).
The objective of this FS is to
address the robustness issue induced by the violation of physical
bounds in the context of heterogeneous approximation techniques.
In this FS we will develop generic limiting techniques that are
robust, efficient, HPC portable, and potentially easy to implement in
LANL's legacy codes.

\vspace*{\baselineskip}

\textit{How:} Enforcing physical bounds is usualy done in the
litearture by limiting slopes or fluxes. Most limiting methods cosnist
of squeezing some accurate (but fragile) approximation towards a
robust (but inaccurate) approximation until some physical bounds are
met. This powerful methodology is unfortunately highly intrusive as it
needs to access and manipulate fluxes and blend high-orderr and
low-order fluxes. The intrusiveness renders this methodology
unapplicable for LANL's legacy codes. It is also extremly curbesome
when deployed on multiphysics problems are the fluxes in this cases
may not be easily indentified.

The key idea that we will deploy in the context of this FS is based on
a algorithm recently proposed by
J.L. Guermond\footnote{J.-L. Guermond,
\url{https://arxiv.org/abs/2606.22839}}.  The method is essentially a
post-processing conservative technique informed of the physics with
some a priori reasonable mathematical knowledge of the bounds that are
expected.  The input for the proposed methods will be a list of
(physical) states a of user-defined concave functionals encoding some
a priori mathematical knowledge on the local bounds tha must be
enforced. The main idea is to locally take from the states that are in
bounds and give to the states that are out of bounds while maintaining
local conservation. The idea sounds simple, but the two main
difficulties that make it non-trivial are that the states and the
budgets are vector-valued and the bounds are often nonliear. The
method has already been demonstrated to be efficient and versatile,
and it is indeed non-intrusive.








\textcolor{red}{\textit{(Suggested section length \(\sim \)1
    page)}}\\ \textcolor{red}{STATE YOUR INNOVATION. COMPARE TO STATE
  OF THE ART. DEFINE SUCCESS AND DETAIL YOUR OUTCOMES (Intellectual
  property? Preliminary data for another proposal? Fill a gap?
  Foundational understanding?)}\\ \textcolor{red}{COMMUNICATE “WHY
  NOW?”}
\vspace*{\baselineskip}

\textcolor{red}{Clearly articulate a technically sound concept that defines the problem well, employs appropriate methods or models, and incorporates testable elements aimed at producing early results or reducing key uncertainties.}
\vspace*{\baselineskip}

\textcolor{red}{Review criteria:}
\begin{quote}
\vspace{-0.3cm}
    \textcolor{red}{\textit{\textbf{Technical Vitality} review criteria weighting 40\%, covering \textbf{Clarity of Problem and Objectives} and \textbf{Innovation and Technical Merit.}}}
\end{quote}

\begin{itemize}[noitemsep,nolistsep]
    \item \textcolor{red}{\textit{What \textbf{technical problem or uncertainty} does this feasibility study address?}}
    \item \textcolor{red}{\textit{Why is resolving this uncertainty important?}}
    \item \textcolor{red}{\textit{What are the \textbf{specific hypotheses, assumptions, or risks} to be tested or reduced?}}
    \item \textcolor{red}{\textit{What is the \textbf{technical design} of the study?}}
    \item \textcolor{red}{\textit{What experiments, models, or analyses will be performed?}}
    \item \textcolor{red}{\textit{How do these tasks directly test the hypotheses or reduce risk?}}
    \item \textcolor{red}{\textit{What \textbf{metrics or criteria} will determine whether hypotheses are supported, disproved, or reframed?}}
    \item \textcolor{red}{\textit{How will the study capture \textbf{useful lessons}, even from negative or inconclusive results?}}
\end{itemize}
\vspace*{\baselineskip}

\textbf{Mission Impact \& Institutional Agility:}


Progress in numerical techniques is crucial for enhancing LANL's
simulation capabilities. Specifically, developing robust forward simulations
that maintain physical states across various multiphysics scenarios will strengthen 
LANL's ability to leverage more efficientlty machine learning and AI tools. 





\textcolor{red}{\textit{(Suggested section length \(\sim \)$\frac{1}{2}$ page)}}\\
\textcolor{red}{Foster cutting edge science, technology, and engineering that ready the Laboratory to respond to current and future mission and national needs.}
\vspace*{\baselineskip}

\textcolor{red}{Review criteria:}\\
\begin{quote}
    \vspace{-0.4cm}\textcolor{red}{\textit{\textbf{Mission Agility} review criteria weighting 20\%, covering \textbf{Mission Agility} and \textbf{Mission Institutional Learning.}}}
\end{quote}

\begin{itemize}[noitemsep,nolistsep]
    \item \textcolor{red}{\textit{What is the potential \textbf{value of the insights} (positive or negative) for:}}
    \begin{itemize}[noitemsep,nolistsep]
        \item \textcolor{red}{\textit{the \textbf{technical community}}}
        \item \textcolor{red}{\textit{LANL's \textbf{mission space}, and}}
        \item \textcolor{red}{\textit{\textbf{future projects} that could build on these findings?}}
    \end{itemize}
    \item \textcolor{red}{\textit{How will the study contribute to \textbf{institutional learning} (roadmaps, datasets, methods, lessons-learned)?}}
\end{itemize}
\vspace*{\baselineskip}

\textbf{Approach, Resources \& Path Forward:}

Ths project will involve the PI (svetlana Tokareva) and two co-PIs who each hold Joint
appointments (JT) with LANL and are  professors of Mathematics at Texas A\&M
University (Jean-Luc Guermond and Matthias Maier). Both J.L. Guermond
and M. Maier are US citizen.

The budget asked for the excecution of the project will divided as
follows: (1) two one-week visits at Texas A\&M by the PI; two one-week
visits at LANL by the two co-PIs during Spring 2027 (two weeks each);
One long visit during the summer (five to six weeks for each
co-PI).

The PI and co-PI are going to focus on the theoretical aspects
of the project in the Fall of 2026 and identify a prototype production
code that will be used to demosntrate the feasability of the proposed
approach. Sping 2027 will be dedicated to testing the theoretical ideas
from demonstrating 





\textcolor{red}{\textit{(Suggested section length \(\sim \)$\frac{3}{4}$ page)}}\\
\textcolor{red}{Articulates a feasibility study that is well conceived, clearly organized, and appropriately scoped, with structure, methods, and milestones designed to produce testable outcomes and generate results that can guide future work.}
\vspace*{\baselineskip}

\textcolor{red}{Review criteria:}
\begin{quote}
\vspace{-0.4cm}
    \textcolor{red}{\textit{\textbf{Research Approach} review criteria total weight 40\%, covering \textbf{Testability and Risk Reduction, Clarity and Organization} and \textbf{Scope, Readiness, Budget, and Duration.}}}
\end{quote}

\begin{itemize}[noitemsep,nolistsep]
    \item \textcolor{red}{\textit{What is the \textbf{plan of work} for the study?}}
    \item \textcolor{red}{\textit{Key \textbf{{tasks, milestones, or deliverables.}}}}
    \item \textcolor{red}{\textit{How will your approach provide \textbf{evidence to validate, refine, or reject} the hypothesis?}}
    \item \textcolor{red}{\textit{What outputs will you produce (data, models, prototypes, analyses)?}}
    \item \textcolor{red}{\textit{What are the \textbf{success criteria}?}}
    \item \textcolor{red}{\textit{How will results - positive or negative - inform \textbf{next steps} (scale up, redesign, or close)?}}
    \item \textcolor{red}{\textit{What are the \textbf{non-technical risks} (facility access, staffing, regulatory, collaboration) and how will you mitigate them?}}
\end{itemize}
\vspace*{\baselineskip}

\textcolor{red}{Proposers should provide enough information for reviewers to assess the scope of work, budget and duration – and that the project is in compliance with \textbf{\href{https://ldrd.lanl.gov/q/public/policy?inline=true}{LDRD Rules and Procedures} (pdf)}}

\begin{itemize}[noitemsep,nolistsep]
    \item \textcolor{red}{\textit{Why are the \textbf{proposed budget} and \textbf{duration} appropriate? (optional bullets: Personnel, Equipment, Materials, Travel)}}
\end{itemize}
\vspace*{\baselineskip}

\textcolor{red}{Since the projects are such short time scales, it should be stated that the work can be turned on immediately if funded. If there is a new activity involved that needs approval, if special materials or lab equipment needs to be procured, if HPC accounts are not in place already, etc. this will raise questions for the reviewers on whether the proposed research can be completed in the time frame provided.}
\vspace*{\baselineskip}

\textcolor{red}{BUDGET RELATED:}\\
\textcolor{red}{If working with external collaborators, communicate timing, any preferred MOUs and academic partners. MAKE SURE YOU COMPLY WITH \textbf{\href{https://ldrd.lanl.gov/q/public/policy?inline=true}{LDRD Rules and Procedures} (pdf)}}\\
\textcolor{red}{FUNDING LIMITS AND PURCHASE DEADLINES.}
\begin{itemize}[noitemsep,nolistsep]
    \item \textcolor{red}{\textbf{Collaborators: External Collaborations are permitted.} No more than 25\% of the project budget may be spent at a collaborator institution \textbf{(7.3.5 Budget Share for Collaborators).}}
    \item \textcolor{red}{\textbf{Purchases:}}
    \begin{itemize}[noitemsep,nolistsep]
        \item \textcolor{red}{M\&S costs must not exceed 25\% of the project’s total budget \textbf{(7.3.2 Materials and Services (M\&S))}}
        \item \textcolor{red}{All single items \> \$5k must be received at least three-months before the project ends. \textbf{(7.4.2 Purchasing Limitations for Final-year Projects)}}
    \end{itemize}
\end{itemize}

\begin{center}
\textit{end of page limited section}
\end{center}
\vspace{-0.7cm}
\noindent\makebox[\linewidth]{\rule{\paperwidth}{0.4pt}}

%%%%%%%%%%%%%%%%% APPENDICES  %%%%%%%%%%%%%%%%%%
\vspace*{\baselineskip}

\textbf{Appendix:} \textcolor{red}{These sections are not page limited.}
\vspace*{\baselineskip}

\textbf{AI Disclosure statement:} \textcolor{red}{\textbf{Substantive contributions by AI must be disclosed.} \\
\textcolor{red}{See \underline{CEA's} "Minimal Use vs. Substantive Use” and "Disclosure Library" resources.}}
\vspace*{\baselineskip}

\textbf{Glossary of Acronyms:} \textcolor{red}{\textbf{Include a list of acronyms (with definitions) used in the proposal.}}
\vspace*{\baselineskip}

%%%%%%%%%%%%%%%%% REFERENCES  %%%%%%%%%%%%%%%%%%
% If submitting a Concept Disruptor, copy this section starting from \renewcommand and ending in \vspace*{\baselineskip} and paste it into your Concept Disruptor template where indicated (look for commented-out References section. 
%%%%%%%%%%%%%%%%%%%%%%%%%%%%%%%%%%%%%%%%%%%%%%%%

\renewcommand{\bibpreamble}{[1.] This is an example reference. NOTE: If you are a LaTeX user, the list should automatically populate when you type ``$\backslash$cite$\{$insert name of .bib reference$\}$" within the BODY of your proposal.}
\bibliographystyle{abbrv}
\bibliography{LDRD_IR_BIB_template}
\textcolor{red}{\textit{(\textbf{Optional.} While not page limited, strongly recommend to not exceed 1 page)}}
\vspace*{\baselineskip}

\textcolor{red}{This page is for references only. DO NOT include comments, notes, etc.  References should not contain images/figures.  Additional content included in the references that are deemed to be proposal-like (e.g. language such as but not limited to “we,” “this proposal,” “our competition”) may provide a substantive advantage and will result in the rejection of the proposal without technical review. Again, be judicious in your citation selection.}
\vspace*{\baselineskip}

\textcolor{red}{Do not include any letters of reference or recommendation in your proposal. They will be discarded without consideration.}
\vspace*{\baselineskip}

\newpage

%%%%%%%%%%%%%%%%%%%%%%%%%%%%%%%%%%%%%%%%%%%%%%
%%%%%%%%% REQUIRED DMSP TEMPLATE  %%%%%%%%%%%%
%%%%%%%%%%%%%%%%%%%%%%%%%%%%%%%%%%%%%%%%%%%%%%

\begin{center}
    \Large \textbf{DMSP for $<$Proposal Title$>$}
    \vspace{0.5cm}
    
    \normalsize \textbf{\textit{\underline{PI:}}} Name (Last, First, Middle Initial); Group; Email
\end{center}
\vspace*{\baselineskip}

\begin{center}
    
\textcolor{red}{Include a coversheet only if your DMSP for your proposal includes Controlled Unclassified Information (CUI). Delete this page if your proposal does not contain CUI.}
\vspace*{\baselineskip}

\textcolor{red}{\textit{Coversheet not included in proposal page count.}}
\vspace*{\baselineskip}

\textcolor{red}{\textbf{Use the coversheet template provided by the LANL Office of Classification:}}\\
\href{https://www.gsa.gov/system/files/SF901-18a.pdf}{SF901, Controlled Unclassified Information Cover Sheet}
\vspace*{\baselineskip}

\textcolor{red}{(Also include the appropriate page markings in the document header as appropriate. For identifying and marking CUI, see P204-1 Controlled Unclassified at int.lanl.gov/policy.}
\vspace*{\baselineskip}

\textcolor{red}{Navigation: int.lanl.gov/policy $>$ Search for Requirements $>$ enter P204-1, and search}
\vspace*{\baselineskip}

\textcolor{red}{Proposers submitting classified content should work with their DC and reference relevant guidance for appropriate markings.)}
\vspace*{\baselineskip}

\textcolor{red}{\textbf{NOTE: CDO and LDRD rely on you, the subject matter expert, to advise CDO and LDRD on the handling of your DMSP document’s Controlled Unclassified Information (CUI). Please reference the Document upload page in the Q system for CUI handling types that are accepted for LDRD proposals.}}

\end{center}

\newpage

\begin{center}
    \Large \textbf{DMSP for $<$Proposal Title$>$}
    \vspace{0.5cm}
\end{center}

\Large{\textcolor{blue}{\textbf{Data Description}}} 
\vspace*{\baselineskip}

\normalsize
\textbf{Project Title:}
\vspace*{\baselineskip}

\textbf{Project Number:}
\vspace*{\baselineskip}

\textbf{Date:}
\vspace*{\baselineskip}

\textbf{Contact Name and Email:}\\
\textit{Who is responsible for management and sharing of this dataset?}\\
\textit{Provide Znumber, Name, Role, ORCID, and email address.}
\vspace*{\baselineskip}

\textbf{Describe your data.}\\
\textit{Briefly describe the data that will be generated or used in this project, including its purpose, content, and context.}\\
\textit{Will the data be updated after initial release? If so, describe the update frequency and versioning approach.}\\
\textit{Example: “The project is expected to generate two types of data (tabular and simulation output) to support [project goal]. The tabular data will be used for... The simulation data will be used for... We plan to update the data after initial release on a yearly basis; specific versioning changing will be documented with the dataset in a README file.”}
\vspace*{\baselineskip}

\textbf{Scientific Domain}\\
\textit{What is the scientific domain of the dataset?}
\vspace*{\baselineskip}

\textbf{Linked Publications or Outputs (Optional)}\\
\textit{Is the dataset in support of or linked to publications or other outputs such as journal articles or technical reports?  If so, please describe (title, DOI, or other information sufficient to identify and locate.)}\\
\textit{Example: “The publicly available data will be sourced from [citation] or None yet; to be added upon publication/deposit.”}
\vspace*{\baselineskip}

\textbf{Landing page (Optional)}\\
\textit{If available, provide the URLs for the dataset’s landing page, repository record, and/or project homepage. (\href{https://int.lanl.gov/org/ddops/aldbus/cio/cdo/_assets/Datasheet Template.docx}{Link to GM-DP Datasheet})}\\
\textit{Example: “None yet; to be added upon publication/deposit.”}
\vspace*{\baselineskip}

\Large{\textcolor{blue}{\textbf{Data types, sources, and standards}}}
\vspace*{\baselineskip}

\normalsize \textbf{Anticipated Data Formats and Modalities}\\
\textit{What data formats, modalities, and aggregation are expected to be generated? For example: file formats (CSV, HDF5, GeoTIFF, etc.), data modality (tabular, image, text, time series, geospatial, simulation output, etc.), aggregation (raw, aggregated, summarized, etc.).}
\vspace*{\baselineskip}

\textit{Example: The project is expected to generate anticipated tabular and simulation data. Formats will likely include CSV or HDF5; modalities include tabular/time‑series; aggregation is expected to be raw and summarized.}
\vspace*{\baselineskip}

\textbf{Sources}\\
\textit{What will the sources of the data be - either generated by the project or external sources of data?  How will the data be generated or collected?}
\vspace*{\baselineskip}

\textit{Example: The project will generate the simulation data using... or the data will be downloaded from the supplementary material of [paper].}
\vspace*{\baselineskip}

\textbf{Quality and Bias}\\
\textit{Are there any known issues with data quality and/or bias in the data?}
\vspace*{\baselineskip}

\textit{Example: Potential modeling bias from boundary conditions; QC steps will be recorded in the data catalog at release. Or unknown at this time.}
\vspace*{\baselineskip}

\textbf{Volume}\\
\textit{What is the anticipated volume / size of the data?}
\vspace*{\baselineskip}

\textit{Example: Order‑of‑magnitude 10–100 GB; final volume will be captured on the data catalog.}
\vspace*{\baselineskip}

\textbf{Timeline}\\
\textit{What is the anticipated timeline of data production?}
\vspace*{\baselineskip}

\textit{Example: Initial data generation targeted for Q3–Q4 FY26.}
\vspace*{\baselineskip}

\textbf{Metadata / Documentation}\\
\textit{What metadata and documentation will accompany the data to support understanding and reuse? (e.g., README files, data dictionaries, metadata schemas, ontologies, etc.). List discipline-specific standards, if used. (\href{https://int.lanl.gov/org/ddops/aldbus/cio/cdo/_assets/Datasheet Template.docx}{Link to GM-DP Datasheet})}
\vspace*{\baselineskip}

\textit{Example: Metadata will be stored along with in the data in CSV files using [domain specific metadata standard] or Metadata will be stored along with in the data in CSV files with key fields [list a few key metadata fields] .... Additional documentation will be stored in README files...}
\vspace*{\baselineskip}

\Large{\textcolor{blue}{\textbf{Related tools, software, and code}}}\\
\normalsize
\textit{List the tools, software, and code that will be required to generate, process, or use the data.}\\
\textit{Include software (with version), equipment, and facilities as applicable.}\\
\textit{Use a delimited list (e.g., semicolon-separated)}
\vspace*{\baselineskip}

\textit{Example: Python 3.x; TensorFlow 2.x; LANL HPC environment; exact versions to be updated at a later date.}
\vspace*{\baselineskip}

\Large{\textcolor{blue}{\textbf{Data Access and Reuse}}}\\
\normalsize
\textbf{Public Access Level}\\
\textit{What is the expected public access level for the data (e.g., unclassified - unlimited release, controlled unclassified information, etc.)?}
\begin{itemize}[noitemsep,nolistsep]
    \item \textit{Unclassified - Unlimited Release}
    \item \textit{CUI - Controlled Unclassified Information}
    \item \textit{Classified}
    \item \textit{Other}
\end{itemize}
\vspace*{\baselineskip}

\textit{Example: Unclassified – Unlimited Release (subject to review).}
\vspace*{\baselineskip}

\textbf{Persistent Identifiers (PIDs)}\\
\textit{Will a Digital Object Identifier (DOI) or other persistent identifier be applied?}
\vspace*{\baselineskip}

\textit{Example: An internal PID will be produced through DIANA. We will also plan to upload the finalized dataset used in the publication to Zenodo which will produce a public DOI.}
\vspace*{\baselineskip}

\Large{\textcolor{blue}{\textbf{Data Security, Preservation, and Sharing}}}
\vspace*{\baselineskip}

\normalsize
\textbf{Data Storage}\\
\textit{Please list where data will be stored during project duration and before review and release. \href{https://docs.google.com/spreadsheets/d/1_TCq0XZ4PuEopOc877CL_8KEhQj63_Jmxn9T_2U1HtQ/edit?usp=sharing}{RL's Generalist Data Repositories and Storage Matrix} - Link to webpage that outlines storage options (LDRD and non-LDRD prices), Q\&A page and CDO POCs}
\vspace*{\baselineskip}

\textit{Example: Short term data storage will be on HPC scratch and campaign space. Long term internal data storage will be through DIANA.}
\vspace*{\baselineskip}

\textbf{Sharing \& Release}\\
\textit{Describe the plan and timeline for data sharing and release, including any embargo periods, access conditions, IP\#’s or updates over time.  Indicate anticipated \href{https://creativecommons.org/share-your-work/cclicenses/}{Creative Commons license} if any.}
\vspace*{\baselineskip}

\textit{Example: Target public release within 6 months post‑project; potential embargo for publications; expected license CC‑BY (final license on data catalog). All publicly released data and code will follow LANL procedures.}
\vspace*{\baselineskip}

\textbf{Data Repositories}\\
\textit{Please list any data repositories that will likely be used to share the data. See NSTC \href{https://rosap.ntl.bts.gov/view/dot/62310}{"Desirable Characteristics of Data Repositories for Federally Funded Research"}. See \href{https://www.re3data.org/}{https://www.re3data.org/} for example repositories.}
\vspace*{\baselineskip}

\textit{Example: We plan to store any publication data on Zenodo.}
\vspace*{\baselineskip}

\textbf{Data Value}\\
\textit{How long do you expect the data to be of value to the scientific community?}
\vspace*{\baselineskip}

\textit{Example: Anticipated community value 5-10 years.}
\vspace*{\baselineskip}

\Large{\textcolor{blue}{\textbf{Provenance \& Motivation}}}\\
\normalsize
\textbf{User Facility/Proposal IDs:}\\
\textit{“If created at a scientific user facility, provide project/proposal identifiers or citation.”}
\vspace*{\baselineskip}

\textit{Example: The dataset will be produced at [user facility]; user facility IDs will be added when assigned.}
\vspace*{\baselineskip}

\textbf{AI Context:}\\
\textit{“Was the dataset created for AI? Describe intended tasks.”}
\vspace*{\baselineskip}

\textit{Example: Yes, the dataset was produced to support [AI task].}
\vspace*{\baselineskip}

\Large{\textcolor{blue}{\textbf{Versioning}}}\\
\normalsize
This table tracks DMSP revisions; dataset versioning will be maintained with the data catalog/repository record.

\begin{table}[h!]
\begin{tabular}{||c|c|c|c||}
\hline
Version & Date & Authors & Changes \\ [1ex] 
 \hline\hline
  &  &  &  \\ 
  \hline
  &  &  &  \\
  \hline
  &  &  &  \\
 \hline
  &  &  &  \\
 \hline
  &  &  &  \\ [1ex] 
 \hline
\end{tabular}
\end{table}

\newpage



%%%%%%%%%%%%% END OF DOCUMENT %%%%%%%%%%%%%%%%%
\end{document}
